% SHARD-ID: CA-97-WHAT-THE-TOWER-REMEMBERS
% SHARD-TITLE: What the Tower Remembers: Rooted AF Invariants and Marked Refinements
% SHARD-SUMMARY: The trailing unital AF limit is classified by its rooted scaled ordered K_0, which can forget the fusion ring, site class and associator even with its canonical trace and stages.
% SHARD-SUMMARY: Fixed-H dyadic shuffles descend with identity K_0 action; abstract diagrams forget H, named finite compressions recover it, and interval/dilation reconstruction remains open.
% SHARD-KEYWORDS: AF algebra, rooted Bratteli diagram, scaled ordered K0, fusion ring, trace, half-braiding, dyadic refinement, marked reconstruction, fidelity, F6
\section{What the Tower Remembers: Rooted AF Invariants and Marked Refinements}\label{sec:what-the-tower-remembers}
% SHARD-KEYWORDS: AF algebra, rooted Bratteli diagram, scaled ordered K0, fusion ring, trace, half-braiding, dyadic refinement, marked reconstruction, fidelity, F6

\subsection{Status, Sources, and Dependencies}

\textbf{Status: WEAKER-RESULT-PROVED} (REVIEW:3--12). Statements 97.1--97.2, 97.4--97.8 and 97.11--97.13 are
\emph{derived}; 97.3, 97.9--97.10 and 97.14 are \emph{negative}; 97.15 is \emph{conditional}; 97.16 is
\emph{conjectural}, an open alternative. No reviewed GAP is promoted. NOTE, REVIEW and INTEGRATION denote files in
\path{explorations/2026-10-25-astra-r3/F6-what-the-tower-remembers/}; all statement locators refer there.
\textbf{CA-88 cells (INTEGRATION:43--62):} (any; \(\mathsf L0\); \(\mathsf R0^u\)) gains classification in
\emph{trailing geometry}; (any; \(\mathsf L0+H\); \(\mathsf R0^u\)) gains fixed-\(H\) dyadic compatibility and marked
extraction; (any; \(\mathsf L0+H\); \(\mathsf R1\)) gains only the canonical tracial zero-drift instance on that schedule.
The \(\mathsf R4\) insufficiency annotation strengthens; no \(\mathsf R2\)--\(\mathsf R5\) attainment, \(\mathsf R0^c\),
\(\mathsf L3\) fidelity or full/periodic cell moves. Bare-\(\mathsf L0\) dyadic existence is not proved.
Dependencies: CA-77 (77.1--77.3), CA-80 (80.8--80.10), CA-88 (grid), CA-90 (90.H, 90.1, 90.3), CA-93 (93.2, 93.8);
CA-71--CA-73 and CA-94 supply boundary comparisons only. CONVENTIONS (a), (b), (r), (am), (aq), (as), (av) apply;
the additional scoped notation below is proposed in SHARD\_META.md. Convention (ar) is unrelated Ward notation.

% Source: explorations/2026-10-25-astra-r3/F6-what-the-tower-remembers/NOTE.md:20-63,67-287,302-508,520-707,811-914 -- definitions and proofs.
% Source: explorations/2026-10-25-astra-r3/F6-what-the-tower-remembers/REVIEW.md:3-12,29-96,100-163,196-214,259-265,291-293 -- binding verdicts and corrections.
% Source: explorations/2026-10-25-astra-r3/F6-what-the-tower-remembers/INTEGRATION.md:18-62,64-114,122-173 -- grades, cells, amendments, conventions, target.
% Source: explorations/2026-10-25-astra-r3/briefs/F6-what-the-tower-remembers.md:82-110 -- original conjecture, corrected here.
% Source: references/text/PenneysUnitaryFusionCategories.md:219-237,277-288,319-338,394-419,470-487,1010-1019 -- splitting, graded categories, coherence, dimensions.
% Source: references/text/EtingofNikshychOstrikFusionCategories.txt:1459-1487,1700-1706 -- unique positive dimension character and unitary spherical dimensions.
% Source: references/text/EtingofNikshychOstrikFusionCategories.md:39-51 -- twisted graded fusion categories.
% Source: references/manifest/SOURCES.md:34-35,55-57 -- primary-source registration.
% Source: literature/md/2010.11121/2010.11121.md:166-180 -- strict norm limits and compatible states.
% Source: report/sections/77_pointed_carrier_invariants.tex:47-86 -- trailing versus corner and fully marked reconstruction.
% Source: report/sections/90_local_unital_lift.tex:59-141,175-211,215-230 -- H, shuffle, finite coherence and converse boundary.
% Source: report/sections/80_lattice_sewing_and_fidelity_certificate.tex:153-159,168-214 -- insufficiency and auxiliary-state distinction.
% Source: report/sections/93_normal_limit_expectations_and_fidelity.tex:57-114,156-198 -- represented expectations and conditional fidelity.
% Source: report/sections/88_layered_targets.tex:84-154 -- target grid; CONVENTIONS.md:631-667,1450-1468,1524-1543,1564-1576,1607-1620.
% Source: report/sections/94_number_changing_collision_dynamics.tex:212-234 -- named reaction data and naturality, not R5.
% Source: explorations/2026-10-25-astra-r3/F6-what-the-tower-remembers/checks/check_f6.py:1-132 -- author invariants.
% Source: explorations/2026-10-25-astra-r3/F6-what-the-tower-remembers/checks/output.txt:1-7 -- author GREEN record.
% Source: explorations/2026-10-25-astra-r3/F6-what-the-tower-remembers/checks/README.md:1-29 -- mutation commands and scope.
% Source: explorations/2026-10-25-astra-r3/F6-what-the-tower-remembers/checks/review_run_output.txt:1-9 -- restored runs and intended failures.
% Source: explorations/2026-10-25-astra-r3/F6-what-the-tower-remembers/checks/review_twisted_refinement.py:1-146 -- independent exact check.
The checks bundle records six exact author groups, two restored review runs and five intended assertion failures
(REVIEW:165--194). These test finite identities, not classification or infinite reconstruction. No local primary AF
classification source was found (NOTE:291--300; REVIEW:241--245); 97.8 uses the reviewed finite-stage proof.

\subsection{Setting and Rooted Invariant}

\(\Cc\) is unitary fusion, \(O\ne0\), \(\mathcal D=\langle O\rangle\), \(I=\Irr\mathcal D\), \(W_0=\one\),
\(W_L=O^{\otimes L}\) left-associated, \(A_L=\operatorname{End}_\Cc(W_L)\), \(i_L(a)=a\otimes1_O\).
Use column incidence \(N[y,x]=\dim\Hom(y,x\otimes O)\), \(m_0=e_\one\), \(m_L=N^Le_\one\),
\(S_L=\{x:m_L[x]>0\}\); absent blocks are omitted. Let \(\kappa_L:A_L\to A_\infty\) be the trailing limit maps.
An AF algebra is the norm completion of an increasing finite-dimensional union. \(K_0\) is the group completion
of stable Murray--von Neumann projection classes; its cone comes from projections and its scale is \([1]\).
\(\Tr\) below is ordinary unnormalised matrix trace, rank-one value one; \(d_x\) is the canonical positive dimension.
No point is needed until stated; then \(\eta:\one\to O\), \(\eta^\dagger\eta=1\), as in (am).

\textbf{Theorem 97.1 [derived; any; \(\mathsf L0\); \(\mathsf R0^u\), trailing] (NOTE:67--153; REVIEW:29--37).}
\textbf{Assume} the setting above. \textbf{Prove} the maps \(i_L\) form a local unital injective tower, with
\[
 A_L\cong\bigoplus_{x\in S_L}M_{m_L[x]}(\mathbb C),\qquad
 (K_0(A_\infty),K_0(A_\infty)^+,[1])=\varinjlim_L(\mathbb Z^{S_L},\mathbb Z_+^{S_L},m_L;N|_{S_{L+1}\times S_L}).
\]
Its Bratteli diagram is the \emph{rooted reachable unfolding}: vertices \(S_L\), edges \(N[y,x]\).
\par\(\langle1\rangle1\). Orthogonal splitting gives
\(\Hom(y,W_L\otimes O)=\bigoplus_x\Hom(x,W_L)\otimes\Hom(y,x\otimes O)\); \(i_L(a)\) acts as \(a_x\otimes1\).
Thus multiplicities are \(N[y,x]\), and \(\sum_yN[y,x]d_y=d_xd_O>0\) prevents any source block disappearing.
Tensor dagger/interchange gives unitality and locality; the maximum block norm proves isometry and the AF completion.
\par\(\langle1\rangle2\). Stable projection classes in a finite sum of matrices are nonnegative block ranks;
the unit has rank vector \(m_L\), and inclusions act by \(N\).
\par\(\langle1\rangle3\). Continuity uses only finite-dimensional stages. Approximate a limit projection \(p\)
by self-adjoint finite-stage \(b\), \(\|b-p\|<\epsilon<1/4\); \(q=1_{[1/2,\infty)}(b)\) has \(\|q-p\|<2\epsilon<1\).
For \(z=qp+(1-q)(1-p)\), \(z-1=(q-p)(2p-1)\), so \(z(z^*z)^{-1/2}\) conjugates \(p\) to \(q\).
Conversely, approximate a limit partial isometry \(v\) between finite projections \(p,q\) by finite \(w\).
For \(x=qwp\), both \(x^*x\) and \(xx^*\) are invertible in their corners; \(x(x^*x)^{-1/2}\) witnesses equivalence there.
\par\(\langle1\rangle4\). Every class and relation is therefore witnessed at a finite stage; group completion gives
the displayed cone and unit. Common-stage coordinate order also proves directedness, unperforation and interpolation. QED.

\textbf{Proposition 97.2 [derived; any; \(\mathsf L0\); \(\mathsf R0^u\), trailing] (NOTE:155--210; REVIEW:38--41).}
Put \(G_N=\varinjlim(\mathbb Z^I,\mathbb Z_+^I;N)\), \(u=[e_\one,0]\). Then
\(K_0(A_\infty)\cong I_u:=\{g\in G_N:-nu\le g\le nu\text{ for some }n\in\mathbb N\}\), with inherited cone and unit \(u\).
If \(N\) is primitive, \(I_u=G_N\); a pointed \(O\) makes \(N\) primitive on \(I\).
\par\(\langle1\rangle1\). Extend supported vectors by zero. Eventual equality and positivity show an order embedding.
Every supported \(v\) is bounded by \(\pm nm_L\); conversely such bounds, witnessed at one later level, force support in \(S_L\).
\par\(\langle1\rangle2\). All simples of \(\mathcal D\) occur in nonnegative powers. Indeed their union \(S\) is
tensor/summand closed. For \(x\in S\), \(P[z,y]=N_{xy}^z d_z/(d_xd_y)\) is column-stochastic with strictly positive
stationary weights proportional to \(d_y^2\), by reciprocity. No edge exits \(S\); stationarity forbids entry, so the
edge \(x^*\to\one\) forces \(x^*\in S\). Thus \(S=I\).
\par\(\langle1\rangle3\). A point makes \(S_L\) increasing, hence eventually full. If \(O^K\) contains every simple,
for each \(x,y\) choose \(z\subset x^*\otimes y\); reciprocity gives \(N_{xz}^y>0\), so \((N^K)[y,x]>0\).
Conversely \(N^p>0\) gives a full supported tail (no zero rows); the limits then share that tail. QED.

\textbf{Negative 97.3 [negative; any; \(\mathsf L0\); \(\mathsf R0\)] (NOTE:212--222; REVIEW:42--43).}
For an invertible generator \(O=g\) of a cyclic group of order \(q>1\), \(A_L=\mathbb C\), \(i_L=\mathrm{id}\),
so rooted \(K_0=\mathbb Z\). But \(N\) is a cyclic permutation and \(G_N=\mathbb Z^q\) with coordinate cone;
under \(N^{-L}\), \(u\) lies on one axis and is not an order unit of \(G_N\). Tensor generation alone is insufficient.
Also \(O=0\) kills \(A_0=\mathbb C\) at the first map and has \(d_O=0\); the nonzero hypothesis is essential (NOTE:38--39).

\subsection{Trace and Classification}

\textbf{Theorem 97.4 [derived; any; \(\mathsf L0\); \(\mathsf R0\) with trace] (NOTE:226--278; REVIEW:44--48).}
With \(\lambda=d_O\), \(\tau_L(a)=\lambda^{-L}\sum_{x\in S_L}d_x\Tr(a_x)\) is a compatible tracial state,
unique if \(N\) is primitive. Rank-one weights are \(d_x/\lambda^L\); block-identity weights are \(m_L[x]d_x/\lambda^L\).
\par\(\langle1\rangle1\). \(d^TN=\lambda d^T\), \(d^Tm_L=\lambda^L\) prove normalization, positivity and compatibility.
\par\(\langle1\rangle2\). Any trace has nonnegative rank-one weights \(w_L=N^Tw_{L+1}\), \(m_L^Tw_L=1\).
If \(N^p>0\), then \(m_{L+p}[y]/\lambda^{L+p}\ge c_y:=\min_x(N^p)[y,x]/(\lambda^pd_x)>0\).
Consequently \(r_L[y]=\lambda^Lw_L[y]/d_y\le1/(c_yd_y)\) eventually.
\par\(\langle1\rangle3\). \(P[x,y]=N[y,x]d_y/(\lambda d_x)\) is row-stochastic, \(r_L=Pr_{L+1}\).
For \(s=|I|\), \(\delta_0=\min(P^p)>0\), subtracting \(\delta_0\) from each entry gives
\(\operatorname{osc}(P^pv)\le(1-s\delta_0)\operatorname{osc}(v)\). Iteration and the bound force \(r_L\) constant;
normalization makes it one. Compatibility fixes earlier levels; norm continuity extends the state. QED.

\textbf{Proposition 97.5 [derived; any, open; \(\mathsf L0+H\); \(\mathsf R1\), tracial only] (NOTE:280--287; REVIEW:111--114).}
On the fixed pointed-\(H\) schedule of 97.11, \(\tau_{2L}\alpha_L=\tau_L\): spherical trace is invariant under the
shuffle conjugation and trailing insertion \(L\) times preserves \(\tau\). Thus every Wilson drift is zero, giving
the compatible limit state on this declared schedule. This selects neither a critical state nor a conformal vacuum.

\textbf{Lemma 97.6 [derived; any; category-free; \(\to\mathsf R0\)] (NOTE:302--319; REVIEW:49--50).}
For \(B=\bigoplus_iM_{b_i}\), \(C=\bigoplus_jM_{c_j}\), all sizes positive, unital \(*\)-maps \(B\to C\) modulo target-unitary
conjugacy correspond exactly to nonnegative integer matrices \(T\) with \(Tb=c\); the induced \(K_0\) map is \(T\).
\par\(\langle1\rangle1\). In each target block, matrix units split the source's central summands into equivalent
orthogonal ranges of dimensions \(T[j,i]\), giving \(a_i\mapsto a_i\otimes1_{T[j,i]}\).
\par\(\langle1\rangle2\). Dimension counting gives \(Tb=c\); these blocks construct every such map. Matching bases
give target-unitary conjugacy, and minimal projection ranks give \(T\). QED.

\textbf{Proposition 97.7 [derived; any; \(\mathsf L0\); \(\mathsf R0^u\), trailing] (NOTE:321--342; REVIEW:51--53).}
A unit-preserving based-ring isomorphism carrying \([O]\) to \([O']\) gives a stage-filtered trailing-tower isomorphism.
\par\(\langle1\rangle1\). It matches \(N,m_L,S_L\). Inductively choose a block-label-preserving isomorphism at the next
stage; its two maps from the preceding stage have equal multiplicities. Conjugate the new target using 97.6.
\par\(\langle1\rangle2\). The resulting commuting isometries extend to the completions. No interval, concatenation,
placement or dyadic mark is promised. Canonical positive dimensions already depend only on the based ring. QED.

\textbf{Theorem 97.8 [derived; any; category-free; \(\to\mathsf R0\)] (NOTE:344--413; REVIEW:54--61).}
For sequential unital AF limits \(A=\varinjlim A_n\), \(B=\varinjlim B_m\) of \emph{finite-dimensional} algebras with
injective connecting maps, \(A\cong_*B\) iff \((K_0(A),K_0(A)^+,[1_A])\cong(K_0(B),K_0(B)^+,[1_B])\).
\par\(\langle1\rangle1\). A \(*\)-isomorphism preserves all projection data. Conversely, lift the finitely many positive
columns of a unit-preserving stage-to-limit group map to one finite stage, by 97.1's continuity proof. Its unit
discrepancy vanishes at a later stage; there \(Tb=c\) exactly. Lemma 97.6 realises this lift as a unital \(*\)-map.
\par\(\langle1\rangle2\). Equality of two stage group maps into a limit is equality of finitely many columns at
some later stage; 97.6 upgrades this to target-unitary conjugacy of their algebra maps.
\par\(\langle1\rangle3\). Alternately lift the group isomorphism and its inverse along increasing cofinal stages.
Using step 2, conjugate each new target to impose exact triangles
\(\psi_k\phi_k=i_{n_k,n_{k+1}}\), \(\phi_{k+1}\psi_k=i'_{m_k,m_{k+1}}\), leaving earlier triangles unchanged.
No stage \(K_0\) map is assumed injective. Each algebra map is injective because its completed triangle is.
\par\(\langle1\rangle4\). The triangles give compatible isometric maps and inverse maps on dense unions, hence
inverse maps on the completions. QED. This classifies the \emph{unmarked limit}, not its filtration or endomorphisms.

\subsection{Information Loss}

\textbf{Negative 97.9 [negative; any; \(\mathsf L0\); \(\mathsf R0\)] (NOTE:422--475; REVIEW:62--66).}
For any finite group \(G\), \(|G|=q\), normalized unitary cocycle \(\omega\), and regular site
\(O=\bigoplus_{g\in G}g\) in \(\Hilb_{\mathrm{fd}}(G,\omega)\), put \(J_q=(1)_{x,y}\). Then
\[\begin{gathered}
 N=J_q,\quad m_L=q^{L-1}\mathbf1\ (L\ge1),\quad (K_0,K_0^+,[1])=(\mathbb Z[1/q],\mathbb Z[1/q]_{\ge0},1),\\
 A_\infty\cong\varinjlim(M_{q^L},b\mapsto b\otimes1_q).
\end{gathered}\]
\par\(\langle1\rangle1\). Exactly one \(g=x^{-1}y\) gives each edge; \([v,L]\mapsto q^{-L}\sum_xv_x\) gives the
ordered group: zero-sum vectors die under \(J_q\), and positive sums become positive vectors; the unit maps to one.
\par\(\langle1\rangle2\). For \(B_L=\bigoplus_gM_{q^{L-1}}\), block-diagonal inclusion into \(M_{q^L}\) and replication
into \(B_{L+1}\) give exact interleaving triangles. This is also the limit for \(\mathrm{Vec},O'=q\one\), with fusion-ring
rank one versus \(q>1\). Both inputs are pointed, generating and primitive. For fixed \(\mathrm{Vec}\), \(q\one\) and
\(q^k\one\), \(q,k\ge2\), have identical limits by cofinal blocking, but different site classes. QED.

\textbf{Negative 97.10 [negative; any; \(\mathsf L0\); \(\mathsf R0\), obstruction to decoding] (NOTE:477--508,668--686; REVIEW:67--69,79--81).}
If two normalized cocycles on \(G\) are unrelated by automorphism and coboundary, their regular-site categories are
monoidally inequivalent but have isomorphic stage-filtered trailing towers with matching canonical trace.
\par\(\langle1\rangle1\). Any tensor equivalence permutes invertible simples by a group isomorphism; normalized
tensorator scalars \(\beta\) force the cocycle ratio, after relabelling, to be
\[(\partial\beta)(a,b,c)=\frac{\beta(b,c)\beta(a,bc)}{\beta(ab,c)\beta(a,b)}.\] The assumed obstruction forbids it.
\par\(\langle1\rangle2\). Apply 97.7 and the identical dimension weights of 97.4. The explicit E-F instance below
realises the obstruction. A decoder depending only on this output record would assign one record two inequivalent
inputs, impossible. QED. No Morita-inequivalence claim is made. Fully marked CA-77 reconstruction can already retain
\(F\) at \(\mathsf R0\); a compatible trace need not restore it. Thus ``\(F,R\) first enter at \(\mathsf R1\)'' is false.

\subsection{What the Refinement Record Adds}

Here \(O\) is pointed and \(H=(t_Y)\) is CA-90's unitary relative half-braiding on the full \(\mathcal D\):
\(t_Y:Y\otimes O\to O\otimes Y\), natural, \(t_\one=1\), \(t_{Y\otimes Z}=(t_Y\otimes1_Z)(1_Y\otimes t_Z)\),
\(t_Y(1_Y\otimes\eta)=\eta\otimes1_Y\). Coherent associators are understood. \(\mathsf L1\) supplies such \(H\);
\(H\) need not supply a braiding on \(\mathcal D\). Set \(\alpha_L=\operatorname{Ad}_{B_L}(\cdot\otimes1_{W_L})\),
where \(B_L\) interleaves old/new strands, crossing old \(j\) with new \(i\) exactly when \(i<j\) (CA-90).
\textbf{Proposition 97.11 [derived; any, open; \(\mathsf L0+H\); \(\mathsf R0^u\)] (NOTE:520--558; REVIEW:70--72,259--265).}
For \emph{one fixed} \(H\) at every length,
\begin{equation}\label{eq:97-square}
 \alpha_{L+1}(a\otimes1_O)=\alpha_L(a)\otimes1_O^{\otimes2},\qquad
 \delta\kappa_L=\kappa_{2L}\alpha_L,\qquad \delta_*=\mathrm{id},\quad\tau\delta=\tau .
\end{equation}
Here \(\delta\) is an injective unital AF endomorphism; its powers form a discrete semigroup.
\par\(\langle1\rangle1\). If \(\Sigma_L:O\otimes W_L\to W_L\otimes O\) moves the last old strand across the first \(L\) new ones,
\(B_{L+1}=(B_L\otimes1_O^{\otimes2})(1_{W_L}\otimes\Sigma_L\otimes1_O)\). Only disjoint crossings are reordered.
\par\(\langle1\rangle2\). The right factor commutes with \(a\) on the first \(L\) factors, proving the square;
the isometric maps therefore define \(\delta\) on the union and its completion.
\par\(\langle1\rangle3\). Inner conjugation fixes projection classes, so \(\delta_*[v,L]=[N^Lv,2L]=[v,L]\).
The trace assertion is 97.5's conjugation calculation. QED; no continuous flow is constructed.
Reviewer correction: \(H\) already implies \(t_{12}t_{23}t_{12}=t_{23}t_{12}t_{23}\), by tensoriality and naturality
for \(t_O\); step 1 needs no extra braid axiom (REVIEW:259--265).

\textbf{Proposition 97.12 [derived; any, open; \(\mathsf L0+H\); \(\mathsf R0^u\)] (NOTE:560--588; REVIEW:73--75).}
The \emph{abstract} dyadic diagram at \(L_n=2^nL_0\), \(L_0\ge1\), is isomorphic to the trailing cofinal diagram.
\par\(\langle1\rangle1\). Put \(j_n(a)=a\otimes1_{W_{L_n}}\), \(U_0=1\), \(U_{n+1}=j_n(U_n)B_{L_n}^\dagger\).
Direct multiplication gives \(\operatorname{Ad}_{U_{n+1}}\alpha_{L_n}=j_n\operatorname{Ad}_{U_n}\). QED.
Its negative consequence is loss of \(H\) and \(F\) discrimination (E-H, E-F); comparison requires \(H\) on both inputs.
These conjugacies need not also preserve trailing maps, interval inclusions or named junctions.

\textbf{Theorem 97.13 [derived; any, open; \(\mathsf L0+H\) and named marks; \(\mathsf R0^u\)] (NOTE:590--634; REVIEW:76--78).}
Retain the pointed tensor record and the following typed placements, adjoints and compression identifications:
\[\begin{gathered}
 r=\eta\otimes1_O,\ l=1_O\otimes\eta:O\to W_2,\quad h=rl^\dagger\in A_2,\\
 i=1_O\otimes1_O\otimes\eta\otimes\eta,\ j=\eta\otimes1_O\otimes1_O\otimes\eta:W_2\to W_4.
\end{gathered}\]
For the supplied-\(H\) formula, \(t_O=j^\dagger\alpha_2(h)i\); the tensor marks determine all of \(H\).
\par\(\langle1\rangle1\). Naturality for \(\eta\) gives \(t_O(\eta\otimes1)=1\otimes\eta\). With \(B=1\otimes t_O\otimes1\),
\(B^\dagger i=1\otimes\eta\otimes1\otimes\eta\), \((h\otimes1\otimes1)B^\dagger i=j\), and \(Bj=jt_O\).
Thus \(j^\dagger B(h\otimes1\otimes1)B^\dagger i=t_O\), since \(j^\dagger j=1\).
\par\(\langle1\rangle2\). Tensoriality fixes \(t_{W_n}\); for each isometric summand \(v:Y\to W_n\),
\(t_Y=(1_O\otimes v^\dagger)t_{W_n}(v\otimes1_O)\). Proposition 97.2 exhausts simples; naturality makes this
independent of splitting, and direct sums finish. QED. This recovers supplied \(H\), not \(H\) from arbitrary local lifts.

\textbf{Negative 97.14 [negative; any, open; \(\mathsf L0\) and lengthwise data; \(\mathsf R0\)] (NOTE:646--651,891--914; REVIEW:88--90).}
Lengthwise finite CA-90 axioms and dyadic composability do not imply \eqref{eq:97-square}. In E-H below,
\(\alpha_2^+(h)\otimes1^{\otimes2}\) has coefficient \(+1\) where \(\alpha_3^-(h\otimes1)\) has \(-1\).
Both maps separately obey the finite axioms. This refutes automatic descent, not the fixed-\(H\) construction.

\textbf{Theorem 97.15 [conditional; chiral/named chiral subnet; \(\mathsf L0+L4\) / \(\mathsf L1+L4\); \(\mathsf R4\) given \(\mathsf R3\)] (NOTE:688--707; REVIEW:82--83,131--137).}
\textbf{Assume} a declared completely rational chiral net \(\Alg\) on a separable Hilbert space and a fully faithful
unitary tensor functor \(\Phi:\Cc\to\mathrm{DHR}_{\mathrm{fin}}(\Alg)\). On its actual four-interval inclusion
\(\mathcal N=\Alg(I_1)\vee\Alg(I_3)\subset\mathcal M=(\Alg(I_2)\vee\Alg(I_4))'\), assume increasing represented
\(\mathcal N_k\subset\mathcal M_k\) boundedly exhaust both algebras, normal expectations \(E_k:\mathcal M_k\to\mathcal N_k\),
and one normal state \(\psi\) on \(\mathcal M\), faithful on \(\mathcal N\), preserved by every \(E_k\). Require
\(b:=\liminf K(E_k)\le\operatorname{Dim}\Cc=\sum_xd_x^2\), with cost \(K\) as in (av).
\textbf{Prove} \(\mu_\Alg=\operatorname{Dim}\Cc\) and \(\Phi\) is onto all finite sectors up to equivalence; if \(\Phi\) is
braided at \(\mathsf L1\), the equivalence is braided and \(\Cc\) modular. BY CA-93 Lemma 93.2 and Theorem 93.8.
This is a conditional restatement, not an application: an AF trace supplies none of these represented certificates.
The full setting needs an identified chiral subnet or CA-93's specified doubled net, not the diagonal extension's DHR category.

\subsection{Effect on other shards}
\emph{Amendment pointers (2026-10-25, round 3; INTEGRATION:64--114).} CA-77 keeps its corner and fully marked
reconstructions distinct from 97.8; CA-71--CA-73 keep their corner qualifications. CA-90 gains 97.11--97.14,
without resolving its converse; 97.11 records the reviewer's correction to 90.3's self-Yang--Baxter parenthetical.
CA-80's insufficiency warning extends by 97.10; CA-93 gains only the boundary in 97.15; CA-88 gains exactly the
cells listed above. CA-94 retains its reaction/tensor marks and its distinction from \(\mathsf R5\). CA-74 is unaffected.

\subsection{What This Does Not Claim}
No classification of interval nets or AF endomorphisms by \(K_0\), no effective dimension-group decision algorithm,
no universal loss of \(F\) at \(\mathsf R0\), and no recovery from arbitrary lifts. No periodic seam, parity covariance,
critical-state selection, modes, conformal limit, central-charge lift, characters or coherent CFT synthesis is constructed.

\subsection{Examples (Instances Only)}
All entries have (any; \(\mathsf L0\) plus the displayed declarations; \(\mathsf R0\)); their negative conclusions are
97.3, 97.9, 97.10, 97.14. Exact data: NOTE:839--914; REVIEW:84--90. These are tests of the general statements.
\begin{center}\footnotesize
\begin{tabular}{p{2.4cm}p{11.7cm}}\toprule
Instance & Exact calibration / counterexample \\\midrule
\raggedright E-P: \(\Hilb_{\mathrm{fd}}(C_3),O=g\) & \(A_L=\mathbb C\); rooted \(K_0=\mathbb Z\), full stationary \(G_N=\mathbb Z^3\), \(u\) on one axis. \\
E-R / E-V & Regular sites of \(C_4\) and \(C_2\times C_2\): \(N=J_4\), \(m_L=4^{L-1}\mathbf1\); only the former has a basis element of order four. Both limits equal that for \(\mathrm{Vec},O'=4\one\), of ring rank one. \\
E-O / E-0 & \(\mathrm{Vec},O=2\one,4\one\): same limit by even-stage blocking, different \([O]\). Any \(\Cc,O=\one\): \(A_\infty=\mathbb C\), \(\mathcal D=\mathrm{Vec}\); the ambient complement is invisible. \\
E-F: twisted \(C_2\), \(O=\one\oplus g\) & Labels \(a,b,c,d\in\{0,1\}\), \(\omega_+=1\), \(\omega_-(a,b,c)=(-1)^{abc}\). Pentagon exponents both reduce to \(bcd+abd+acd+abc\) mod 2; normalized \(\partial\beta(1,1,1)=1\), whereas \(\omega_-(1,1,1)=-1\), with no nontrivial automorphism. Both give the same traced \(J_2\) tower. Both admit \(H\): crossing scalar \(r_+=1\); \(r_-=1\) except \(r_-(1,1)=i\), whose only nontrivial tensor identity is \(1=-i^2\). \\
\raggedright E-H: untwisted \(C_2\), \(O=\one\oplus g\) & \(t^\pm(a,b)=(\pm1)^{ab}\,\mathrm{swap}\). Bilinearity gives \(H\); a pointed block-scalar gauge cannot change the \(g,g\) sign. The \((g,g,\one,\one)\to(\one,g,g,\one)\) entry of \(\alpha_2^\pm(h)\) is \(\pm1\). Appending two vacancies gives the same \(+1/-1\) mismatch for \(\alpha_2^+\otimes1^2\) versus \(\alpha_3^-\) in 97.14. \\\bottomrule
\end{tabular}
\end{center}

\subsection{Single Next Target: Intermediate-Record Faithfulness}

\textbf{Target 97.16 [conjectural; any, open; \(\mathsf L0+H\); \(\mathsf R0\) information only] (NOTE:636--644,811--827; REVIEW:196--214).}
For arbitrary pointed tensor-generating inputs \((\mathcal D,O,\eta,H)\), retain exactly
\(\mathcal R=(A_\infty,\{A_{[j,k]}\}_{1\le j\le k},\delta)\), where
\(A_{[j,k]}=\kappa_k(1_{W_{j-1}}\otimes\operatorname{End}_\Cc(W_{k-j+1}))\), extending (as).
Record isomorphisms must carry each named interval algebra onto its counterpart and intertwine \(\delta\);
no point, corner, \(h,i,j\) or concatenation marks remain. Prove that record isomorphism forces a unitary tensor
equivalence transporting \(\eta,H\), by recovering 97.13's missing marks up to that gauge, or give an explicit
counterexample preserving \emph{all} these constraints simultaneously. Neither 97.12 nor 97.13 answers this question.
